\providecommand{\Res}[2]{\qty[round-mode=figures,round-precision=3]{\fpeval{#1}}{#2}}
\providecommand{\Val}[1]{\num[round-mode=figures,round-precision=4]{\fpeval{#1}}}
\providecommand{\ResSci}[2]{\qty[round-mode=figures,round-precision=3,exponent-mode=scientific]{\fpeval{#1}}{#2}}

\subsubsection*{Question a) Force d'une charge adjacente}
\begin{enumerate}[label=\arabic*)]
    \item \textbf{Variables connues et recherchées :} $q = \qty{\v{q}e-6}{C}$, $a = \qty{\vsi{a}}{m}$, $k = \qty{8.99e9}{N.m^2/C^2}$ ; variable recherchée : $F_1$
    \item \textbf{Équations pertinentes :} \[ F_1 = \frac{k\,q^2}{a^2} \]
    \item \textbf{Isolation de la variable recherchée :} l'équation est déjà isolée.
    \item \textbf{Résolution numérique :} \[ F_1 = \frac{\num{8.99e9} \cdot (\num{\v{q}e-6})^2}{(\num{\vsi{a}})^2} = \Res{8.99e9*\vsi{q}^2/\vsi{a}^2}{N} \]
    \item \textbf{Réponse :} $F_1 = \mathbf{\Res{8.99e9*\vsi{q}^2/\vsi{a}^2}{N}}$
\end{enumerate}

\subsubsection*{Question b) Force de la charge en diagonale}
\begin{enumerate}[label=\arabic*)]
    \item \textbf{Variables connues et recherchées :} $q$, distance diagonale $a\sqrt{2}$ ; variable recherchée : $F_2$
    \item \textbf{Équations pertinentes :} \[ F_2 = \frac{k\,q^2}{(a\sqrt{2})^2} = \frac{k\,q^2}{2a^2} \]
    \item \textbf{Isolation de la variable recherchée :} l'équation est déjà isolée.
    \item \textbf{Résolution numérique :} \[ F_2 = \frac{\num{8.99e9} \cdot (\num{\v{q}e-6})^2}{2 \cdot (\num{\vsi{a}})^2} = \Res{8.99e9*\vsi{q}^2/(2*\vsi{a}^2)}{N} \]
    \item \textbf{Réponse :} $F_2 = \mathbf{\Res{8.99e9*\vsi{q}^2/(2*\vsi{a}^2)}{N}}$
\end{enumerate}

\subsubsection*{Question c) Force résultante}
\begin{enumerate}[label=\arabic*)]
    \item \textbf{Variables connues et recherchées :} $F_1 = \Res{8.99e9*(\v{q}e-6)^2/(\vsi{a})^2}{N}$ (deux fois, à angle droit), $F_2 = \Res{8.99e9*(\v{q}e-6)^2/(2*(\vsi{a})^2)}{N}$ ; variable recherchée : $F$
    \item \textbf{Équations pertinentes :} les deux forces adjacentes, perpendiculaires, donnent $\sqrt{2}\,F_1$ le long de la diagonale, dans le même sens que $F_2$ :\[ F = \sqrt{2}\,F_1 + F_2 \]
    \item \textbf{Isolation de la variable recherchée :} \[ F = \frac{k q^2}{a^2}\left(\sqrt{2} + \tfrac{1}{2}\right) \]
    \item \textbf{Résolution numérique :} \[ F = \sqrt{2} \cdot \Val{8.99e9*(\v{q}e-6)^2/(\vsi{a})^2} + \Val{8.99e9*(\v{q}e-6)^2/(2*(\vsi{a})^2)} = \Res{8.99e9*\vsi{q}^2/\vsi{a}^2*(sqrt(2) + 0.5)}{N} \]
    \item \textbf{Réponse :} $F = \mathbf{\Res{8.99e9*\vsi{q}^2/\vsi{a}^2*(sqrt(2) + 0.5)}{N}}$, le long de la diagonale, vers l'extérieur du carré (répulsion).
\end{enumerate}
